\documentclass{article}
\usepackage[a4paper, total={8in, 10in}]{geometry}
\usepackage{amsmath}

\title{Find the two nearest points on the two line segments}
\author{Kostic Srecko}

\begin{document}
\maketitle

\section*{Solution}

Let the following be line segments bounding points:

\[
\mathbf{A}_0 = \begin{pmatrix} x_{a0} \\ y_{a0} \\ z_{a0} \end{pmatrix}, \quad
\mathbf{A}_1 = \begin{pmatrix} x_{a1} \\ y_{a1} \\ z_{a1} \end{pmatrix}, \quad
\mathbf{B}_0 = \begin{pmatrix} x_{b0} \\ y_{b0} \\ z_{b0} \end{pmatrix}, \quad
\mathbf{B}_1 = \begin{pmatrix} x_{b1} \\ y_{b1} \\ z_{b1} \end{pmatrix}
\]

Let $\mathbf{u}$ and $\mathbf{v}$ be the direction vectors for two segments $\mathbf{L}_A$ and $\mathbf{L}_B$, respectively:

\begin{align*}
    \mathbf{u} &= \mathbf{A}_1 - \mathbf{A}_0 \\
    \mathbf{v} &= \mathbf{B}_1 - \mathbf{B}_0
\end{align*}

Let the following be the two line segments $\mathbf{L}_A$ and $\mathbf{L}_B$:

\begin{align*}
    \mathbf{L}_A(t) &= \mathbf{A}_0 + t\mathbf{u} \\
    \mathbf{L}_B(s) &= \mathbf{B}_0 + s\mathbf{v}
\end{align*}

Let $\mathbf{w}$ be a direction vector function representing the distance between two points on the line segments:

\[
\mathbf{w}(t, s) = \mathbf{L}_A(t) - \mathbf{L}_B(s)
\]

The shortest line between two line segments is perpendicular to both.

\[
\begin{cases}
    \mathbf{w}(t, s) \cdot \mathbf{u} = 0 \\
    \mathbf{w}(t, s) \cdot \mathbf{v} = 0
\end{cases}
\]

Expand $\mathbf{w}(t, s)$:

\begin{align*}
    \mathbf{w}(t, s) &= \mathbf{L}_A(t) - \mathbf{L}_B(s) \\
    \mathbf{w}(t, s) &= (\mathbf{A}_0 + t\mathbf{u}) - (\mathbf{B}_0 + s\mathbf{v}) \\
    \mathbf{w}(t, s) &= \mathbf{A}_0 + t\mathbf{u} - \mathbf{B}_0 - s\mathbf{v} \\
    \mathbf{w}(t, s) &= \mathbf{A}_0 - \mathbf{B}_0 + t\mathbf{u} - s\mathbf{v}
\end{align*}

Isolate variables:

\begin{align*}
    \mathbf{w}(t, s) \cdot \mathbf{u} &= (\mathbf{A}_0 - \mathbf{B}_0) \cdot \mathbf{u} + t(\mathbf{u} \cdot \mathbf{u}) - s(\mathbf{v} \cdot \mathbf{u}) \\
    \mathbf{w}(t, s) \cdot \mathbf{v} &= (\mathbf{A}_0 - \mathbf{B}_0) \cdot \mathbf{v} + t(\mathbf{u} \cdot \mathbf{v}) - s(\mathbf{v} \cdot \mathbf{v})
\end{align*}

To simplify, isolate variable $\mathbf{w}_0$:

\[
\mathbf{w}_0 = \mathbf{A}_0 - \mathbf{B}_0
\]

Substitute back in the system, and isolate constant to the other side:

\begin{align*}
    \mathbf{w}_0 \cdot \mathbf{u} + t(\mathbf{u} \cdot \mathbf{u}) - s(\mathbf{v} \cdot \mathbf{u}) = 0 \\
    \mathbf{w}_0 \cdot \mathbf{v} + t(\mathbf{u} \cdot \mathbf{v}) - s(\mathbf{v} \cdot \mathbf{v}) = 0 \\
    \\
    t(\mathbf{u} \cdot \mathbf{u}) - s(\mathbf{v} \cdot \mathbf{u}) = - \mathbf{w}_0 \cdot \mathbf{u} \\
    t(\mathbf{u} \cdot \mathbf{v}) - s(\mathbf{v} \cdot \mathbf{v}) = - \mathbf{w}_0 \cdot \mathbf{v} \\
\end{align*}

To simplify, isolate scalars to variables:

\[
a = \mathbf{u} \cdot \mathbf{u}, \quad
b = \mathbf{u} \cdot \mathbf{v}, \quad
c = \mathbf{v} \cdot \mathbf{v}, \quad
d = \mathbf{u} \cdot \mathbf{w}_0, \quad
e = \mathbf{v} \cdot \mathbf{w}_0
\]

Substitute constants:

\[
\begin{cases}
    at - bs = -d \\
    bt - cs = -e
\end{cases}
\]

Solve the system using Cramer's rule:

\begin{align*}
    D &= b^2 - ac \\
    D_t &= dc - eb \\
    D_s &= bd - ae \\
    t &= \frac{D_t}{D} \\
    s &= \frac{D_s}{D} \\
    t &= \frac{dc - eb}{b^2 - ac} \\
    s &= \frac{bd - ae}{b^2 - ac}
\end{align*}

Let $\mathbf{P}$ and $\mathbf{Q}$ be the two closest points:

\begin{align*}
    \mathbf{P} &= \mathbf{L}_A(t) \\
    \mathbf{Q} &= \mathbf{L}_B(s)
\end{align*}

However, this finds $\mathbf{P}$ and $\mathbf{Q}$ on an infinite line segments $\mathbf{L}_A$ and $\mathbf{L}_B$ respectively. To solve that problem, we check first if $t$ is outside of bounds, and compute $s$.

\subsection*{Derive the closest point equation's parameter}

This section derives the two closest points in terms of the line equation parameter.

Before all, if the lines are parallel, then there is no intersection.

\begin{align*}
    D \approx 0
\end{align*}

First, we start by clamping the computed $s$ from the Cramer's rule:

\begin{align*}
    s_{\text{init}} &= \text{clamp}(\frac{D_s}{D}, 0, 1)
\end{align*}

\textit{Note: It does not matter which of $\frac{D_t}{D}$ or $\frac{D_s}{D}$ you start
from --- the double projection applies regardless, as long as you know which parameter
belongs to which segment.}

Below are the steps and explanations to derive the point to line projection, using the formula, written down in terms of already computed variables.

$A$ - Point to project onto the line

$P$ - Starting point of the line

$m$ - Direction vector of the line

The following computes the point along the infinite line:

\begin{align*}
    k &= \frac{(A - P) \cdot m}{m \cdot m} \\
    Q &= A + km \\
    Q & = A + \frac{(A - P) \cdot m}{m \cdot m} m
\end{align*}

The computed $s_{\text{init}}$ lies on $B_0$, $B_1$ segment.

The vector $v$ is the direction vector of $B$ segment.

The segment direction vector $m$ is: $m = u$

Also $b = v \cdot u$, and $d = w_o \cdot u$

\begin{align*}
    Q &= A + \frac{(P - A) \cdot m}{m \cdot m} \cdot m \\
    P &= B_0 + s_{\text{init}} v  \\
    A &= A_0 \\
    P - A &= B_0 + s_{\text{init}} v - A_0 \\
    P - A &= -(A_0 - B_0) + s_{\text{init}} v \\
    w_0 &= A_0 - B_0 \\
    P - A &= - w_0 + s_{\text{init}} v \\
    P - A &= s_{\text{init}} v - w_0 \\
    m &= u \\
    Q &= A_0 + \frac{(s_{\text{init}} v - w_0) \cdot m}{m \cdot m} \cdot m \\
    Q &= A_0 + \frac{(s_{\text{init}} v - w_0) \cdot u}{u \cdot u} \cdot u \\
    Q &= A_0 + \frac{s_{\text{init}} v \cdot u - w_0 \cdot u}{u \cdot u} \cdot u \\
    b &= v \cdot u \\
    d &= w_o \cdot u  \\
    a &= u \cdot u \\
    Q &= A_0 + \frac{s_{\text{init}} (v \cdot u) - (w_0 \cdot u)}{(u \cdot u)} \cdot u \\
    Q &= A_0 + \frac{s_{\text{init}} b - d}{a} \cdot u \\
    k &= \text{clamp}(\frac{s_{\text{init}} b - d}{a}, 0, 1) \\
    Q &= A_0 + k \cdot u \\
\end{align*}

We clamp the $k$. Then, we project back using computed $k$, and we are certainly within the two segments, between the two points on the shortest distance from each other.

Note: The logic should be reusable in terms of input parameters, which returns a scalar. The scalar represents clamped parametric value.

Then, we repeat the process, project onto another line, and extract the parameter.

\begin{align*}
    Q &= A + \frac{(P - A) \cdot m}{m \cdot m} \cdot m \\
    m &= v \\
    P &= A_0 + ku \\
    A &= B_0 \\
    P - A &= A_0 + ku - B_0 \\
    P - A &= A_0 - B_0 + ku \\
    w_0 &= A_0 - B_0 \\
    P - A &= w_0 + ku \\
    Q &= B_0 + \frac{(w_0 + ku) \cdot v}{v \cdot v} \cdot v \\
    Q &= B_0 + \frac{w_0 \cdot v + k (u \cdot v)}{v \cdot v} \cdot v \\
    b &= u \cdot v \\
    c &= v \cdot v \\
    e &= w_0 \cdot v \\
    Q &= B_0 + \frac{kb + e}{c} \cdot v \\
    l &= \text{clamp}(\frac{kb + e}{c}, 0, 1) \\
    Q &= B_0 + l \cdot v \\
\end{align*}

At the very end, if the distance between the two closest points is below some $\epsilon$ threshold, we conclude that there is an intersection.

\begin{align*}
    \text{distance}(W, Q) \approx 0
\end{align*}

\end{document}
